\chapter{Sum-product partition inference}
\label{ch:sumproduct}
\label{app:marginals}
\chapterpromise{Forward/backward identities, evidence, segment-count uncertainty, boundary marginals, and model-averaged Bayes curves.}

\label{sec:dp}

For $k\in\{0,\dots,k_{\max}\}$ and $j\in\{0,\dots,n\}$, let $L_{k j}$ denote the evidence of
the prefix $y_{1:j}$ under a segmentation into $k$ segments (after integrating out segment
parameters), and let $R_{k i}$ denote the evidence of the suffix $y_{i+1:n}$ under $k$
segments. We adopt the conventions
\[
L_{0,0}=1,\quad L_{0,j}=0\ (j>0),\qquad
R_{0,n}=1,\quad R_{0,i}=0\ (i<n).
\]

\subsection{Forward/backward recursions}
Given single-block evidences $A^{(0)}_{ij}$, the forward and backward recursions are
\begin{equation}
L_{k+1,j}
 = \sum_{h=k}^{j-1} L_{k,h}\,A^{(0)}_{h j},\qquad
R_{k+1,i}
 = \sum_{h=i+1}^{n-k} A^{(0)}_{i h}\,R_{k,h},
\label{eq:LR}
\end{equation}
for $k\ge0$, $1\le j\le n$, and $0\le i\le n-1$, where sums are restricted to valid indices.

\begin{proposition}[Forward--backward duality and total-evidence identity]
\label{prop:fb-duality}
Under Assumption~\ref{ass:standing-offline} and Assumption~\ref{ass:factorizing-prior}, with
prior-adjusted block scores $\widetilde A^{(0)}_{ij}=A^{(0)}_{ij}\gamma_x(i,j)$
\eqref{eq:Atilde}, let $\{\widetilde L_{k,j}\}$ and $\{\widetilde R_{k,i}\}$ denote the prefix and
suffix sum-product messages of \eqref{eq:LR}. Then for every admissible
$k\in\{1,\dots,k_{\max}\}$:
\begin{enumerate}[label=\textnormal{(\roman*)}]
\item \emph{Boundary conditions.}
$\widetilde L_{0,0}=1$, $\widetilde L_{0,j}=0$ for $j>0$, $\widetilde R_{0,n}=1$, and
$\widetilde R_{0,i}=0$ for $i<n$. Furthermore $\widetilde L_{k,j}=0$ whenever $j<k$ and
$\widetilde R_{k,i}=0$ whenever $i>n-k$.
\item \emph{Total-evidence identity.}
$\widetilde L_{k,n}=\widetilde R_{k,0}=\sum_{t\in\mathcal T_{k,n}^{+}}\prod_{q=1}^k \widetilde A^{(0)}_{t_{q-1}t_q}$,
i.e.\ the forward-evidence terminal and the backward-evidence initial value both equal the same
prior-adjusted partition sum.
\item \emph{Boundary marginal product form.}
For interior boundary position $h\in\{p,\dots,n-(k-p)\}$ and $1\le p\le k-1$,
\[
p(t_p=h\mid y,k)=\frac{\widetilde L_{p,h}\,\widetilde R_{k-p,h}}{\widetilde L_{k,n}},
\]
and the boundary-event marginal of Definition~\ref{def:posterior-summaries}(a) satisfies
$\mathrm{MB}_h(k)=\sum_{p=1}^{k-1}p(t_p=h\mid y,k)$.
\end{enumerate}
\end{proposition}

\begin{proof}
Boundary conditions follow from definition; the support constraints $j\ge k$ and $i\le n-k$
restate that a $k$-segment prefix requires at least $k$ observations and a $k$-segment suffix requires at least $k$ observations.
For (ii), iterate \eqref{eq:LR} forward from $\widetilde L_{0,0}$: by induction on $k$,
$\widetilde L_{k,j}=\sum_{t_0<\cdots<t_k=j}\prod_{q=1}^k\widetilde A^{(0)}_{t_{q-1}t_q}$,
so $\widetilde L_{k,n}$ enumerates all $k$-segmentations of $\{1,\dots,n\}$. The same induction
backward from $\widetilde R_{0,n}$ gives $\widetilde R_{k,i}=\sum_{i=t_0<\cdots<t_k=n}\prod_{q=1}^k\widetilde A^{(0)}_{t_{q-1}t_q}$,
and $\widetilde R_{k,0}$ enumerates the same set. Both equal the asserted partition sum.
For (iii), split each segmentation at the $p$th boundary: a segmentation with $t_p=h$ is uniquely
described by its first $p$ blocks (counted by $\widetilde L_{p,h}$) and its last $k-p$ blocks
(counted by $\widetilde R_{k-p,h}$). Their product, divided by $\widetilde L_{k,n}$, is the
conditional probability that $t_p=h$. Summing over $p\in\{1,\dots,k-1\}$ yields the boundary-event
marginal because the boundaries are strictly ordered, so for fixed $h$ at most one $p$ satisfies
$t_p=h$ in any given $t$.
\end{proof}

The identity $\widetilde L_{k,n}=\widetilde R_{k,0}$ is also one of the DP invariants used in the
runtime correctness checklist of Section~\ref{sec:alg-checklist}; Proposition~\ref{prop:fb-duality}
states the algebraic basis of that numerical check.

\paragraph{Implementation.}
Algorithm~\ref{alg:dp-core} presents the algebraic sum-product recursions corresponding to
\eqref{eq:LR}. It is a mathematical specification rather than the recommended floating-point
implementation: Section~\ref{sec:alg-dp-impl} gives the implementable log-space version using
\texttt{logsumexp}. The algebraic routine takes a triangular block-evidence array and returns the
forward messages, backward messages, and segment-count posteriors needed for boundary and moment
extraction.

\begin{algorithm}[H]
\caption{Generic algebraic DP for EF segmentation (prefix/suffix form)}
\label{alg:dp-core}
\KwIn{Block evidences $\widetilde{A}^{(0)}_{ij}$ for $0\le i<j\le n$; $k_{\max}$; $C_k$ (default $C_k=\binom{n-1}{k-1}$)}
Initialize $\widetilde{L}_{0,0}=1$, $\widetilde{L}_{0,j>0}=0$; $\widetilde{R}_{0,n}=1$, $\widetilde{R}_{0,i<n}=0$\;
\For{$k=0$ \KwTo $k_{\max}-1$}{
  \For{$j=1$ \KwTo $n$}{
    $\widetilde{L}_{k+1,j}\leftarrow \sum_{h=k}^{j-1} \widetilde{L}_{k,h}\widetilde{A}^{(0)}_{h j}$\;
  }
  \For{$i=0$ \KwTo $n-1$}{
    $\widetilde{R}_{k+1,i}\leftarrow \sum_{h=i+1}^{n-k} \widetilde{A}^{(0)}_{i h} \widetilde{R}_{k,h}$\;
  }
}
Compute $P(k\mid y)\propto p(k)\,\widetilde{L}_{k n}/C_k$ (cf.\ \eqref{eq:post-k}) and choose $\widehat{k}$\;
\For{$p=1$ \KwTo $\widehat{k}-1$}{
  For $h\in\{p,\dots,n-(\widehat{k}-p)\}$ set $P(t_p=h\mid y,\widehat{k})\leftarrow \widetilde{L}_{p,h}\widetilde{R}_{\widehat{k}-p,h}/\widetilde{L}_{\widehat{k}n}$\;
  Record the marginal mode $\arg\max_h P(t_p=h\mid y,\widehat{k})$ as a diagnostic only\;
}
\KwOut{$\{\widetilde{L}_{k j}\}$, $\{\widetilde{R}_{k i}\}$, $P(k\mid y)$, $\widehat{k}$, and boundary-marginal diagnostics}
\end{algorithm}

\subsection{Incorporating segment-factorized boundary priors}
\label{subsec:priorinc}

Suppose the boundary prior given $k$ factorizes over segment lengths as in \eqref{eq:partition-prior-factor}:
\[
p(t\mid k) = \frac{1}{C_k}\prod_{q=1}^k g(\Delta_x(t_{q-1},t_q)),
\qquad
C_k := \sum_{t:\,0=t_0<\cdots<t_k=n}\ \prod_{q=1}^k g(\Delta_x(t_{q-1},t_q)).
\]
Then we can absorb $g$ into block evidences by defining
\begin{equation}
\widetilde{A}^{(0)}_{ij} := A^{(0)}_{ij}\,\gamma_x(i,j),
\qquad
\widetilde{A}^{(r)}_{ij} := A^{(r)}_{ij}\,\gamma_x(i,j).
\label{eq:Atilde}
\end{equation}
Running the DP \eqref{eq:LR} with $\widetilde{A}^{(0)}_{ij}$ computes $\widetilde{L}_{k n}$,
the \emph{unnormalized} evidence sum over all segmentations weighted by $g$. The normalizer
$C_k$ can be computed by the same DP with $A^{(0)}_{ij}\equiv 1$:
\begin{equation}
C_k
 = \sum_{t}\prod_{q=1}^k g(\Delta_x(t_{q-1},t_q))
 = L^{(g)}_{k n},
\quad\text{where}\quad
L^{(g)}_{k+1,j}=\sum_{h=k}^{j-1}L^{(g)}_{k,h}\,\gamma_x(h,j),\ L^{(g)}_{0,0}=1.
\label{eq:Ck-general}
\end{equation}
In the index-uniform case $g\equiv 1$, \eqref{eq:Ck-general} reduces to $C_k=\binom{n-1}{k-1}$.

\subsection[Posteriors over segment counts, boundary marginals, and joint MAP segmentations]{Posteriors over $k$, boundary marginals, and joint MAP segmentations}
Let $p(k)$ be a prior on segment count. Under index-uniform $p(t\mid k)$ we have
$C_k=\binom{n-1}{k-1}$ and
\[
P(y\mid k) = \frac{L_{k n}}{C_k},
\qquad
P(k\mid y)\propto p(k)\,\frac{L_{k n}}{C_k}.
\]
For general segment-factorized $p(t\mid k)$, replace $L$ by $\widetilde{L}$ and $C_k$ by
\eqref{eq:Ck-general}. In all cases we can write compactly
\begin{equation}
P(k\mid y)\ \propto\ p(k)\,\frac{\widetilde{L}_{k n}}{C_k},
\qquad k=1,\dots,k_{\max}.
\label{eq:post-k}
\end{equation}
Let $\widehat{k}=\arg\max_k P(k\mid y)$ denote the MAP segment count.

Conditional on $k$, the marginal posterior for the $p$th boundary $t_p$ is
\begin{equation}
P(t_p=h\mid y,k)
 = \frac{\widetilde{L}_{p,h}\,\widetilde{R}_{k-p,h}}{\widetilde{L}_{k n}},
\qquad
h\in\{p,\dots,n-(k-p)\},
\label{eq:boundary-post}
\end{equation}
where $(\widetilde{L},\widetilde{R})$ are computed from $\widetilde{A}^{(0)}$ as in
\eqref{eq:LR}. The denominator is the unnormalized evidence under the length-weighted partition
prior; the normalizer $C_k$ cancels in \eqref{eq:boundary-post} because the same $C_k$ multiplies
every boundary vector in the numerator and denominator. (When $g\equiv 1$, tildes can be dropped.)

\begin{remark}[Boundary-event marginal]
\label{rem:boundary-event}
A related object we use repeatedly is the \emph{boundary-event marginal}
\begin{equation}
P(b_i=1\mid y,k)
 := \sum_{p=1}^{k-1} P(t_p=i\mid y,k)
 = \sum_{p=1}^{k-1}\frac{\widetilde{L}_{p,i}\,\widetilde{R}_{k-p,i}}{\widetilde{L}_{k n}},
 \qquad 1\le i\le n-1,
\label{eq:boundary-event}
\end{equation}
the posterior probability that interior index $i$ coincides with \emph{some} boundary, without
conditioning on which boundary. The objects $P(t_p=h\mid y,k)$ and $P(b_i=1\mid y,k)$ are
distinct: the former is a probability distribution over $h$ for each fixed $p$ (and sums to $1$ over
$h$), while the latter is a per-index Bernoulli probability summing to $k-1$ over $i$. Calibration
experiments in Section~\ref{sec:experiments} target $P(b_i=1\mid y,k)$; a full derivation is given
in Appendix~\ref{app:marginals}.
\end{remark}

A useful diagnostic is the \emph{marginal boundary mode}
\[
\widetilde{t}^{\,\mathrm{marg}}_p(k)
  \in \arg\max_h P(t_p=h\mid y,k).
\]
However, maximizing the marginals independently need not yield an ordered boundary vector and, in
general, does \emph{not} recover the joint MAP segmentation.

\begin{remark}[A small counterexample]
\label{rem:marg-vs-joint}
Consider $n=5$ and $k=3$, so the two interior boundaries satisfy $1\le t_1<t_2\le4$.
For example, let the joint posterior put masses $0.30,0.28,0.22,0.20$ on the four ordered
configurations $(1,4),(2,3),(2,4),(3,4)$, respectively. The marginal modes are then
$t_1=2$ and $t_2=4$, so the independently selected marginal-mode vector is $(2,4)$, even though
the joint MAP configuration is $(1,4)$. This is exactly the kind of discrepancy the max-sum DP
resolves.
\end{remark}

The joint MAP segmentation is the solution of
\begin{equation}
\widehat{t}^{\,\mathrm{MAP}}(k)
 \in \arg\max_{0=t_0<\cdots<t_k=n}
 \Bigg\{\sum_{q=1}^k \log \widetilde{A}^{(0)}_{t_{q-1} t_q}\Bigg\},
\label{eq:joint-map-k}
\end{equation}
with the convention that any $k$-specific prior terms such as $\log p(k)-\log C_k$ are added only
when selecting among different values of $k$. Equation~\eqref{eq:joint-map-k} is solved exactly by
a max-sum DP with backpointers. Inline, the max-sum recursion reads
\begin{equation}
M_{k+1,j} = \max_{h\in\{k,\dots,j-1\}}\big\{M_{k,h} + \log \widetilde{A}^{(0)}_{hj}\big\},
\qquad
h^\star(k+1,j) \in \arg\max_{h} \big\{M_{k,h} + \log \widetilde{A}^{(0)}_{hj}\big\},
\label{eq:max-sum-inline}
\end{equation}
with initial conditions $M_{0,0}=0$, $M_{0,j}=-\infty$ for $j>0$, and backtracking from
$(k,n)$ as detailed in Section~\ref{sec:algorithms}.
In the sequel, whenever we report an \emph{exported MAP segmentation}, we mean the joint optimizer
$\widehat{t}^{\,\mathrm{MAP}}:=\widehat{t}^{\,\mathrm{MAP}}(\widehat{k})$ rather than a vector of marginal
boundary modes.

\subsection{Segment-level and regression-curve posterior moments}
For segment-wise posterior moments attached to a \emph{fixed} segmentation
$t=(t_0,\dots,t_k)$, let segment $q\in\{1,\dots,k\}$ correspond to block
$(i,j]=(t_{q-1},t_q]$. Then
\begin{equation}
\mu^{(r)}_q(t)
 = \frac{\widetilde{A}^{(r)}_{ij}}{\widetilde{A}^{(0)}_{ij}}
 = \frac{A^{(r)}_{ij}}{A^{(0)}_{ij}},
\qquad r=1,2,\dots,
\label{eq:segmom-seg}
\end{equation}
where the final equality holds because the prior factor $\gamma_x(i,j)$ cancels. In particular, for
an exported joint MAP segmentation $\widehat{t}^{\,\mathrm{MAP}}$, the segment summaries are obtained by
setting $t=\widehat{t}^{\,\mathrm{MAP}}$ in \eqref{eq:segmom-seg}.

To compute the Bayes regression curve within a fixed $k$ (typically $k=\widehat{k}$),
define, for each block $(i,j]$ and moment order $r$,
\[
F^{(r)}_{ij}(k)
 := \frac{1}{\widetilde{L}_{k n}}\sum_{m=1}^k
       \widetilde{L}_{m-1,i}\,\widetilde{A}^{(r)}_{ij}\,\widetilde{R}_{k-m,j}.
\]
Then the $r$th posterior moment of the latent mean at index $t$ given $k$ is obtained by summing
$F^{(r)}_{ij}(k)$ over all candidate blocks $(i,j]$ that contain the index $t$ (i.e., all blocks
with $i<t\le j$), so that each candidate segmentation contributes exactly once through the block
that actually covers $t$:
\begin{equation}
\widehat{\mu}'^{(r)}_t
 := \mathbb{E}[(\mu'_t)^r\mid y,k]
 = \sum_{i<t\le j} F^{(r)}_{ij}(k),
\qquad t=1,\dots,n.
\label{eq:segmom}
\end{equation}
In particular, $\widehat{\mu}'_t := \widehat{\mu}'^{(1)}_t$ is the Bayes regression curve and
$\mathbb{V}[\mu'_t\mid y,k]
 = \widehat{\mu}'^{(2)}_t - (\widehat{\mu}'^{(1)}_t)^2$.
If a fully $k$-marginalized curve is desired, compute the same conditional curve for each admissible
$k$ and average the resulting moments using $P(k\mid y)$ from \eqref{eq:post-k}. The empirical
sections report curves conditional on the selected $\widehat k$ unless explicitly stated otherwise.

\begin{proposition}[Block-covering decomposition for Bayes-curve moments]
\label{prop:block-covering-decomposition}
Fix $k$ and an index $t\in\{1,\dots,n\}$. Under the support and factorization hypotheses of
Theorem~\ref{thm:dp-correctness}, the conditional posterior moment of the latent mean at $t$ is
obtained by summing the block contributions $F^{(r)}_{ij}(k)$ over exactly those candidate blocks
$(i,j]$ with $i<t\le j$. No segmentation contributes more than once to the sum in
\eqref{eq:segmom}.
\end{proposition}

\begin{proof}
Every admissible ordered boundary vector partitions $\{1,\dots,n\}$ into disjoint blocks, hence
contains a unique block $(i,j]$ satisfying $i<t\le j$. Conditioning on that covering block splits
the remaining boundaries into an $(m-1)$-segmentation of the prefix $y_{1:i}$ and a
$(k-m)$-segmentation of the suffix $y_{j+1:n}$ for some segment position $m$. The prefix, block, and
suffix contributions multiply as
$\widetilde L_{m-1,i}\,\widetilde A^{(r)}_{ij}\,\widetilde R_{k-m,j}$, and summing over $m$ and
over all covering blocks gives \eqref{eq:segmom}. Disjointness of the covering block proves that no
segmentation is double counted.
\end{proof}

\begin{theorem}[Correctness of the DP]
\label{thm:dp-correctness}
Assume finite block evidences on all admissible blocks, a finite segment-count range, and a
segment-factorized conditional partition prior satisfying Assumption~\ref{ass:factorizing-prior}
with the boundary-coordinate convention of Section~\ref{sec:notation}. Under the EF--conjugate
single-segment model of Section~\ref{sec:ef}, or under any other block model whose exact evidence
array is supplied, the recursions \eqref{eq:LR} compute exact marginal evidences for all
prefixes/suffixes:
$\widetilde{L}_{k j}$ and $\widetilde{R}_{k i}$ are exact evidence sums over all
$k$-segmentations of the corresponding prefix/suffix. Moreover:
\begin{enumerate}
\item The posterior $P(k\mid y)$ given by \eqref{eq:post-k} is exact.
\item The boundary marginal $P(t_p=h\mid y,k)$ given by \eqref{eq:boundary-post} is exact.
\item The segment-wise moments \eqref{eq:segmom-seg} and regression-curve moments \eqref{eq:segmom}
are exact posterior moments conditional on $k$.
\end{enumerate}
The separate max-sum recursion and backtracking theorem for the joint MAP partition is stated and
proved immediately after its statement in Theorem~\ref{thm:map-correctness}.
\end{theorem}

\begin{proof}
We argue for general $\widetilde{A}^{(0)}_{ij}$; $g\equiv 1$ is the untilded specialization.

\paragraph{Step 1: $\widetilde{L}_{k j}$ is the prefix evidence sum.}
For any $k$-segmentation $(t_0,\dots,t_k)$ of $y_{1:j}$ with $t_0=0$, $t_k=j$, integrating out
segment parameters yields $p(y_{1:j}\mid t,k)=\prod_{q}\widetilde{A}^{(0)}_{t_{q-1},t_q}$.
Summing over admissible $t$ gives
$\widetilde{L}_{k j}=\sum_{t:\,t_k=j}\prod_{q=1}^k \widetilde{A}^{(0)}_{t_{q-1},t_q}$.

\paragraph{Step 2: forward recursion.}
Conditioning a $(k{+}1)$-segmentation of $y_{1:j}$ on its last boundary $h\in\{k,\dots,j-1\}$
splits it into a $k$-segmentation of $y_{1:h}$ and the final block $(h,j]$; summing over $h$
gives the forward recursion in \eqref{eq:LR}. The backward recursion follows by the same
argument applied to suffixes.

\paragraph{Step 3: posterior over $k$.}
$P(y\mid k)=\widetilde{L}_{k n}/C_k$ by definition of $C_k$; Bayes' rule yields \eqref{eq:post-k}.

\paragraph{Step 4: boundary posteriors.}
Conditioning on $t_p=h$ splits the sequence into a $p$-segmentation of $y_{1:h}$ and a
$(k{-}p)$-segmentation of $y_{h+1:n}$, contributing
$\widetilde{L}_{p,h}\,\widetilde{R}_{k-p,h}$; normalizing by $\widetilde{L}_{k n}$ yields
\eqref{eq:boundary-post}.

\paragraph{Step 5: moments.}
Theorem~\ref{thm:ef-integral} gives $A^{(r)}_{ij}=A^{(0)}_{ij}\,\mathbb{E}[m_\star(\theta)^r\mid(i,j]]$;
the length factor $\gamma_x(i,j)$ cancels in the ratio, so \eqref{eq:segmom-seg} holds.
Equation~\eqref{eq:segmom} follows by grouping segmentation contributions through the unique
block $(i,j]$ that contains index $t$.
\end{proof}

\paragraph{Posterior-evidence sanity checks.}
After every DP run, verify four invariants: normalized segment-count weights
sum to one; the full-sequence forward and backward evidences agree, $\widetilde L_{kn}=\widetilde
R_{k0}$; boundary-event marginals satisfy $\sum_{i=1}^{n-1}P(b_i=1\mid y,k)=k-1$; and the
max-sum backtracking score equals the stored terminal max-sum value. In log-space arithmetic these
checks are evaluated with a declared tolerance, typically a small multiple of machine epsilon times
the number of accumulated summands. A failure of the first two checks is treated as a normalization
or admissibility mismatch; a failure of the boundary-sum check indicates an indexing or support
error; and a failure of the backtracking check indicates inconsistent tie-breaking or stale
backpointers.

\begin{theorem}[Time and space complexity]
\label{thm:dp-complexity}
Assume single-block evidences $A^{(0)}_{ij}$ are available for all $0\le i<j\le n$.
Computing all forward and backward messages via \eqref{eq:LR} requires
$\mathcal{O}(k_{\max}n^2)$ arithmetic operations. The triangular block-evidence matrix requires
$\mathcal{O}(n^2)$ memory, while all forward and backward layers require
$\mathcal{O}(k_{\max}n)$ memory. Max-sum values and backpointers each require
$\mathcal{O}(k_{\max}n)$ additional memory. Boundary-event probabilities for every $k$ and candidate
index require $\mathcal{O}(k_{\max}n)$ output storage. Ordered-boundary distributions for every
triple $(k,p,h)$ require $\mathcal{O}(k_{\max}^2n)$ output storage, or $\mathcal{O}(kn)$ for one
fixed segment count $k$. For $R$ requested reporting moments, retaining all block moment numerators
requires $\mathcal{O}(Rn^2)$ memory and the final Bayes curves require $\mathcal{O}(Rn)$; explicitly
retaining every block-cover contribution over all $k$ can require
$\mathcal{O}(k_{\max}n^2)$ additional workspace but is not necessary when those contributions are
accumulated sequentially.
\end{theorem}

\begin{proof}
For each fixed $k$, the forward recursion computes $L_{k+1,j}$ for $j=1,\dots,n$, and each
$L_{k+1,j}$ sums at most $j-k$ terms. Thus the forward work per $k$ is
$\sum_{j=1}^n \mathcal{O}(j)=\mathcal{O}(n^2)$. The backward recursion is analogous. Repeating
for $k=0,\dots,k_{\max}-1$ gives $\mathcal{O}(k_{\max}n^2)$ time. A message or backpointer array
has one entry per $(k,j)$ state. Boundary-event output has one value per $(k,h)$, whereas the full
ordered-boundary output has one value per $(k,p,h)$ and therefore
$\sum_{k\le k_{\max}}\mathcal O(kn)=\mathcal O(k_{\max}^2n)$ entries. The remaining statements
follow by counting the triangular block arrays, reporting-moment arrays, and optional block-cover
workspace.
\end{proof}

\begin{corollary}[Decision rules for full-vector and coordinatewise losses]
\label{cor:map-01}
For the loss $\ell_{01}(t,t')=\mathbbm{1}\{t\neq t'\}$ on the full boundary vector, the Bayes
estimator is the joint MAP segmentation $\widehat{t}^{\,\mathrm{MAP}}(k)$ from
\eqref{eq:joint-map-k}. For a separable coordinatewise loss and an unconstrained Cartesian action
space, each coordinate can instead be chosen by minimizing its marginal posterior risk; under
coordinatewise 0--1 loss this gives the marginal modes. If the reported action must belong to the
strictly ordered partition space $\mathcal T^+_{k,n}$, coordinatewise marginal modes can be
inadmissible, and the Bayes action must minimize posterior expected loss over that constrained set.
It can differ from both the coordinatewise marginal modes and the joint MAP partition.
\end{corollary}

\begin{proof}
The posterior risk of $t'$ under full-vector 0--1 loss is $1-p(t=t'\mid y,k)$, so any joint
posterior mode minimizes it. Under a separable loss on an unconstrained Cartesian action space, the
posterior risk is a sum of coordinatewise risks, each minimized using the corresponding marginal
posterior. Restricting the action to strictly ordered boundary vectors couples the coordinates, so
independent coordinatewise minimization no longer solves the constrained decision problem.
\end{proof}

\begin{remark}[Memory/time trade-offs for large $n$]
\label{rem:mem-time}
Keeping only the current and previous $k$-layers reduces message memory to $\mathcal{O}(n)$, in
addition to any stored block array, when only terminal fixed-$k$ evidences are required. Posterior
marginals and MAP paths then require retained layers or a documented recomputation pass. Checkpoint
spacing can trade retained layers against recomputation, but its exact cost depends on which
summaries are requested and is therefore an implementation choice rather than a universal bound.
When the $n^2$ block-evidence matrix is prohibitive and every segment score is available in constant
time from prefix summaries, scores can instead be evaluated on demand during the DP sweep. This
uses less block storage but performs up to $\mathcal{O}(k_{\max}n^2)$ segment evaluations because
the same segment can be revisited at different DP layers.
\end{remark}

\begin{remark}[Evidence across segment counts need not be monotone]
\label{prop:mono-k}
The fixed-$k$ evidence $P(y\mid k)=\widetilde L_{kn}/C_k$ is not monotone in general. Adding a boundary changes both the set of admissible partitions and
the product of block evidences, while the normalized conditional prior contributes the normalizer
$C_k$. Consequently, diagnose segment-count behavior through the normalized evidences
$P(y\mid k)$ and posterior weights $P(k\mid y)$ rather than through an unnormalized pathwise
``add one boundary'' heuristic.
\end{remark}

\subsection{End-to-end BayesBreak inference routine}\label{sec:dp-alg}
Collecting the ingredients from Sections~\ref{sec:ef} and \ref{sec:dp}---block-evidence precomputation,
exact sum-product DP over partitions, and max-sum backtracking for exported segmentations---yields an
end-to-end routine suitable for direct implementation. Algorithm~\ref{alg:bayesbreak} summarizes
this workflow in a modular way: the only model-specific input is a function that evaluates (or
approximates) block evidences $\log \mathcal{L}_{ij}$ and, when needed, block moment numerators.

\begin{algorithm}[H]
\caption{BayesBreak$(y,x,w,\texttt{model},\texttt{hyper},k_{\max},p(k),g)$}
\label{alg:bayesbreak}
\KwIn{
  Observations $y_{1:n}$ at locations $x_{1:n}$ (optional); observation-model weights/exposures $w_{1:n}$ (default determined by the chosen likelihood, often $w_i\equiv 1$);\\
  model family $\texttt{model}\in\{\texttt{Gaussian},\texttt{Poisson},\texttt{Binomial},\texttt{BetaObs},\dots\}$; \\
  hyperparameters $\texttt{hyper}$; maximum segment count $k_{\max}$; prior $p(k)$; length factor $g(\cdot)$ (default $g\equiv 1$).
}
\KwOut{
  Posterior $P(k\mid y)$; MAP segment count $\widehat{k}$; boundary marginals $P(t_p\mid y,\widehat{k})$;\\
  exported joint MAP segmentation $\widehat{t}^{\,\mathrm{MAP}}$; segment-level moments; Bayes regression curve; (optional) evidence $P(y)$.
}
\BlankLine
\textbf{1. Precompute block evidences and moments}\;
\For{$0\le i<j\le n$}{
  Compute block summaries $(S_{ij},W_{ij},H_{ij})$ using prefix sums (or running sums)\;
  Call a family-specific routine to obtain $A^{(0)}_{ij},A^{(1)}_{ij},A^{(2)}_{ij}$\;
  Set $\widetilde{A}^{(r)}_{ij}\leftarrow A^{(r)}_{ij}\,\gamma_x(i,j)$ for $r\in\{0,1,2\}$ (cf.\ \eqref{eq:Atilde})\;
}
Compute $C_k$ under the same admissibility and length-prior convention used in the score array; for $g\equiv1$, use $C_k=\binom{n-1}{k-1}$, and otherwise use \eqref{eq:Ck-general} (DP with $A^{(0)}_{ij}\equiv 1$)\;
\BlankLine
\textbf{2. Run the sum-product dynamic program}\;
Compute $\widetilde{L}_{k j},\widetilde{R}_{k i}$ via \eqref{eq:LR} for $k=0,\dots,k_{\max}$ using $\widetilde{A}^{(0)}_{ij}$\;
Compute $P(k\mid y)\propto p(k)\,\widetilde{L}_{k n}/C_k$ via \eqref{eq:post-k}; set $\widehat{k}\leftarrow \arg\max_k P(k\mid y)$\;
\BlankLine
\textbf{3. Boundary marginals}\;
\For{$p=1$ \KwTo $\widehat{k}-1$}{
  For all $h\in\{p,\dots,n-(\widehat{k}-p)\}$, set
  $P(t_p=h\mid y,\widehat{k})\leftarrow \widetilde{L}_{p,h}\widetilde{R}_{\widehat{k}-p,h}/\widetilde{L}_{\widehat{k}n}$ (cf.\ \eqref{eq:boundary-post})\;
}
\BlankLine
\textbf{4. Export a joint MAP segmentation}\;
Run the max-sum/backpointer recursion of Section~\ref{sec:algorithms} using $\log \widetilde{A}^{(0)}_{ij}$ and the chosen count $\widehat{k}$\;
Backtrack to obtain $\widehat{t}^{\,\mathrm{MAP}}=(0=\widehat{t}^{\,\mathrm{MAP}}_0<\cdots<\widehat{t}^{\,\mathrm{MAP}}_{\widehat{k}}=n)$\;
\BlankLine
\textbf{5. Segment-level posterior moments on the exported segmentation}\;
\For{$q=1$ \KwTo $\widehat{k}$}{
  Let $(i,j]\leftarrow (\widehat{t}^{\,\mathrm{MAP}}_{q-1},\widehat{t}^{\,\mathrm{MAP}}_{q}]$\;
  For $r\in\{1,2\}$ set $\widehat{\mu}^{(r)}_q \leftarrow \widetilde{A}^{(r)}_{ij}/\widetilde{A}^{(0)}_{ij}$ (cf.\ \eqref{eq:segmom-seg})\;
}
\BlankLine
\textbf{6. Bayes regression curve within $\widehat{k}$}\;
Using $\widetilde{L},\widetilde{R},\widetilde{A}^{(r)}$ and \eqref{eq:segmom}, compute $\widehat{\mu}'_i$ and $\mathbb{V}[\mu'_i\mid y,\widehat{k}]$ for $i=1,\dots,n$\;
\BlankLine
\textbf{7. Optional: pooling/groups}\;
If multiple subjects or groups are present, pool block evidences (Theorem~\ref{thm:multisubject}) and/or run the finite latent-group template procedure of Chapter~\ref{ch:templates}, then repeat Steps 2--6.
\end{algorithm}


\begin{figure}[H]
\centering
\resizebox{0.96\textwidth}{!}{
\begin{tikzpicture}[node distance=7mm and 12mm]
\node[secondarybox,text width=33mm] (score) {Same local log scores\\$\ell_{ij}$};
\node[primarybox,above right=of score,text width=36mm] (sp) {$\oplus=\logsumexp$\\sum-product};
\node[primarybox,below right=of score,text width=36mm] (ms) {$\oplus=\max$\\max-sum};
\node[successbox,right=of sp,text width=35mm] (marg) {Marginal posterior quantities};
\node[successbox,right=of ms,text width=35mm] (path) {One joint maximizing path};
\draw[arrPrimary] (score)--(sp);\draw[arrPrimary] (score)--(ms);
\draw[arrPrimary] (sp)--(marg);\draw[arrPrimary] (ms)--(path);
\draw[Warning,densely dashed,line width=.8pt] (marg.south)--node[right,font=\scriptsize,text=Warning]{not generally equal}(path.north);
\end{tikzpicture}
}
\caption{Same local scores, different semirings. Sum-product produces marginal posterior quantities; max-sum produces one joint maximizing path.}
\label{fig:semiring-routes}
\end{figure}
